\section{Optimal positive closing offset and its realization}\label{sup:positive}

Here positive offset means $a=R(0)>0$.
The constructions below use the finer offset scale $a=\eta h^3$, with
$\alpha=\eta\mu^2>0$ and $\alpha\to0$; the limiting cubic region
$\mathcal C_B$ has $\alpha\le0$.
Two positive roots precede the second threshold and two persist on the
secondary scale.
The realization below adds the outer negative and positive roots to these
two clusters; the bound of six then excludes further roots in $I$.

\begin{proposition}[Positive-offset six-cycle returns]
\label{corner:positive-offset-realization}
\label{thm:positive-offset-six}
Assume~\cite[(\ref*{main-eq:rank})]{main2026} and use the standard
unfolding on $I$.
Put
\[
 \eta_*=\frac{4}{27B^2},\qquad d_*=\frac{1}{6B}.
\]
For every compact
$\mathcal K\Subset(0,\eta_*)\times(0,d_*)$, there are
$\mu_0,\sigma_0>0$ such that the parameters
\begin{equation}\label{corner:positive-offset-weights}
 m=\mu,\qquad h=\mu^2\sigma,\qquad b=-h,\qquad
 \delta=\mu^3\sigma^2d,\qquad a=\eta h^3
\end{equation}
give exactly one negative and five positive simple fixed points in $I$
for every
\[
 0<\mu<\mu_0,\qquad0<\sigma<\sigma_0,\qquad
 (\eta,d)\in\mathcal K.
\]

Let $u_-(\eta)<u_+(\eta)$ be the two positive roots of
$\eta-z^2+Bz^3$, and put
\[
 r_\pm(d)=\frac{1\pm\sqrt{1-6Bd}}2.
\]
The six states satisfy, uniformly on $\mathcal K$,
\begin{align}
 q_N&=-\frac hC(1+o(1)),\notag\\
 q_{F,\pm}&=h^2\bigl(u_\pm(\eta)^2+o(1)\bigr),\notag\\
 q_{S,\pm}&=\mu^2\sigma^2\bigl(r_\pm(d)^2+o(1)\bigr),\notag\\
 q_O&=\frac{\mu^2}{C^2}(1+o(1)).
                                  \label{corner:positive-offset-roots}
\end{align}
Here $o(1)\to0$ as $\mu+\sigma\to0$, with the positive scales
independent. The constants $\eta_*$ in
\eqref{corner:sharp-positive-offset-limsup} and $d_*$ in
\eqref{corner:nonnegative-offset-bound} are sharp: both are approached
simultaneously by a sequence of these positive-offset returns.
\end{proposition}

\begin{proof}
Uniform convergence of \eqref{corner:positive-offset-weights} to zero
on $\mathcal K$ lets the standard unfolding realize them on one
product of scale intervals.

Below the second threshold $a_\delta=\Gamma_1^{-1}(\delta)$, the passage
and closing expansions give
\begin{equation}\label{corner:positive-first-only-identity}
 D(r^2)=a-hr^2+kA_1(\lambda)r^3+r^4T(r,\lambda),
 \qquad k=1-h,
\end{equation}
where $T$ is jointly smooth on a common signed-root domain.
For $r=hz$ on a compact positive $z$ interval,
\[
 \frac{\Gamma_1(h^2z^2)}{\delta}
   =\frac{\mu z^2}{d}\{1+O(h)\}.
\]
Thus only the first passage is active at these states for all sufficiently
small $\mu$,
uniformly on $\mathcal K$ and independently of $\sigma$. Their
normalized displacement satisfies
\begin{equation}\label{corner:positive-first-only-limit}
 \frac{D(h^2z^2)}{h^3}
 =\eta-z^2+kA_1(\lambda)z^3+hz^4T(hz,\lambda)
 =\eta-z^2+Bz^3+O_{C^1}(\mu).
\end{equation}
Indeed the standard unfolding gives $A_1(\lambda)-B=O(\mu)$ on a bounded
rectangle of scales, and the smooth root formula controls the displayed
remainder and its $z$ derivative.

The limiting polynomial is positive at zero and at infinity.
Its positive critical point at $2/(3B)$ is a minimum, with value
$\eta-\eta_*<0$.
It therefore has two positive simple roots $u_\pm(\eta)$.

The secondary weights \eqref{corner:weights} now have
$\alpha=\eta\mu^2>0$. Substitution in the joint quotient
\eqref{corner:joint-limit} gives, on compact positive $z$ intervals,
\begin{equation}\label{corner:positive-secondary-limit}
 \frac{D(\mu^2\sigma^2z^2)}{\mu^4\sigma^3}
 =z\bigl(z^2-z+\tfrac32Bd\bigr)+O_{C^1}(\mu+\sigma).
\end{equation}
For $0<d<d_*$, its two nonzero limiting roots $r_\pm(d)>0$ have
derivatives $r_\pm(2r_\pm-1)\ne0$. Both passages are active on their
isolating intervals because
\[
 \frac{a_\delta}{\mu^2\sigma^2z^2}
   =\frac{\mu d}{z^2}\{1+o(1)\}\longrightarrow0
\]
uniformly on their isolating intervals.

Finally apply the outer charts
\eqref{corner:negative-large-chart} and
\eqref{corner:positive-large-chart} with $\alpha=\eta\mu^2$.
Their $C^1$ limits $X-X^2$ and $X^{3/2}-CX^2$ have simple roots at
$1$ and $1/C^2$.

All four joint $C^1$ chart limits have compact isolating intervals with uniform
derivative margins on $\mathcal K$. Simple-root continuation gives the
six roots in \eqref{corner:positive-offset-roots} for common
thresholds $\mu_0,\sigma_0$, without a relation between the positive
scales. The windows where only the first passage is active have size
$h^2=\mu^4\sigma^2$, the secondary windows size $\mu^2\sigma^2$,
and the positive outer window size $\mu^2$. Their successive ratios
are $O(\mu^2)$ and $O(\sigma^2)$; the negative window is separated
by sign. The second threshold lies strictly between the first-passage and
secondary windows. All six windows lie in $I$ for small scales,
and~\cite[Proposition~\ref*{main-prop:common-count}]{main2026} excludes
additional fixed points there. The charts also give
\eqref{corner:positive-offset-roots}.

For sharpness choose interior pairs $(\eta_j,d_j)$ tending to
$(\eta_*,d_*)$. Apply the result to each singleton pair and choose
positive $\mu_j,\sigma_j<1/j$ below its two existence thresholds.
The resulting returns have exactly
\[
 \frac{a_j}{h_j^3}=\eta_j\longrightarrow\eta_*,
 \qquad
 \frac{m_j\delta_j}{h_j^2}=d_j\longrightarrow d_*.
\]
\end{proof}

\begin{corollary}[Six-cycle regions with positive offset]
\label{corner:open-six-cycle-region}
\label{cor:six-open}
By~\cite[Corollary~\ref*{main-cor:six-open}]{main2026}, in every admissible
family the parameters giving six simple fixed points in $I$ form an open
set. Under~\cite[(\ref*{main-eq:rank})]{main2026}, this set contains open
subsets with $a>0$, one negative fixed point and five positive fixed points,
accumulating at the reference.
\end{corollary}

\begin{proof}
The openness assertion
is~\cite[Corollary~\ref*{main-cor:six-open}]{main2026}.
For the returns of Proposition~\ref{corner:positive-offset-realization},
shrink the neighbourhood to preserve $a>0$ and the signs of the six
simple roots. An interior parameter pair with both scales tending to zero
gives accumulation at the reference.
\end{proof}
